\documentclass{ximera}

% MATH -----------------------------------------------------------
\newcommand{\nats}{\mathbb N}
\newcommand{\ints}{\mathbb Z}
\newcommand{\rats}{\mathbb Q}
\newcommand{\reals}{\mathbb R}
\newcommand{\complex}{\mathbb C}
\newcommand{\powerset}{\mathscr P}

%-------------------------------------------------------------

\pgfplotsset{compat=1.18}
\begin{document}
\begin{abstract}
 \end{abstract}

    \title{Exemplar Examples 8}
    
    \maketitle

A $500$-liter tank has $100$-liters of pure water when a process begins that pumps in saltwater with $30$ grams per liter of salt at $4$ liters per minute, while the liquid in the tank is kept thoroughly mixed and allowed to drain at $2$ liters per minute.  The process automatically stops when the tank reaches capacity.


\begin{enumerate}

\item Let's have $y(t)$ be the grams of salt in the tank, at time $t\geq 0$ minutes into the process.   Let's derive an IVP that $y(t)$ must satisfy.    


% y'=120-2y/(100+2t), y(0)=0

\vfill

\item Let's solve this IVP.

\vfill
\vfill


% y'+1/(50+t) y=120

% [(50+t)y]'= 120(50+t)=6000+120t ---> (50+t)y=6000t-60t^2 (using initial condition)


% y=60t(100+t)/(50+t)


\item Let's determine the amount of salt in the tank at the time of overflow, and the concentration in grams per liter.


% y=60*200(100+200)/(50+200)=12*1200=14400.  14.4 kg in 500 liters  so 144/5\approx 28.8 grams per liter

\vfill

\end{enumerate}



\newpage


A financial planner is modeling a client's portfolio of financial assets.  They assume that the long-run rate of return will always be at least equivalent to a constant annualized rate of return of $r$ \%, that the initial principle is $A_0$, and that there will be a fixed annual withdrawal rate of $N+\alpha t$ where $N,\alpha$ are fixed positive constants.

\begin{enumerate}
\item Let $y(t)$ be the model's worst case value for the portfolio after $t\geq 0$ years.  What IVP must $y(t)$ satisfy?

% y'=(r/100)y-N-\alpha t, y(0)=A_0

\vfill

\item Let's determine the general solution.

% y'-(r/100)y=-N-\alpha t

% IF:  e^(-(r/100)t)
% e^(-(r/100)t)y'-(r/100)e^(-(r/100)t)y=(-N-\alpha t)e^(-(r/100)t)
% [e^(-(r/100)t)y]'=(-N-\alpha t)e^(-(r/100)t)
% e^(-(r/100)t)y=(N+\alpha t)/(r/100) e^(-(r/100)t) +\alpha/(r/100)^2 e^(-(r/100)t)+C
% y=(N+\alpha t)/(r/100)  +\alpha/(r/100)^2 +Ce^((r/100)t)

% A_0 = N/(r/100)  +\alpha/(r/100)^2 +C
% C=A_0 - N/(r/100)  -\alpha/(r/100)^2 

% y=(N+\alpha t)/(r/100)  +\alpha/(r/100)^2 +(A_0 - N/(r/100)  -\alpha/(r/100)^2)e^((r/100)t)

% (-N-\alpha t) | e^(-(r/100)t)
% -\alpha       | -1/(r/100) e^(-(r/100)t)
% 0             | 1/(r/100)^2 e^(-(r/100)t)

\vfill
\vfill

\item What must $A_0$ satisfy to guarantee an exponential growth term in $y$?

% A_0 > N/(r/100)  +\alpha/(r/100)^2 = (N(r/100)+\alpha)/(r/100)^2

\vfill


\end{enumerate}


\end{document}